
%% Language ---
\usepackage[english]{babel}


%% Settings of siunitx ---
\sisetup{% 
  locale          = UK,
}% end siunitx setup.


\begin{document}


Full Name:
\raisebox{-2pt}{\rule[-.4ex]{5cm}{.4pt}}\hfill
Date: 
\raisebox{-2pt}{\rule[-.4ex]{3cm}{.4pt}}
\hspace*{.5cm}% add some margin space.


% Section Header ---
\vspace{6mm}
% switch '\ifprintanswers' defined in the class file:
%\textbf{\large Einheitenabfrage, Physik \ifprintanswers \hfill \textcolor{AccentColor}{Antworten}% deutsch; es folgt die englische Zeile:
\textbf{\large SI Unit Test \ifprintanswers \hfill \textcolor{AccentColor}{Answers}
\fi
%\hfill 5 Punkte
}% end textbf.
\vspace{3mm}


For each question give one answer only.



%% Fill up to 20 questions overall: --------
%% This copy-paste removes a need for another package (for-loop).
%% The counter 'exerciseID' is from the class file of "exam". 
%% Control the number of questions with this IF here:
%%                    |
%%                    V
\ifnum\theexerciseID<10 {\input{build/Q1}}  \fi% 1
\ifnum\theexerciseID<10 {\input{build/Q2}}  \fi% 2
\ifnum\theexerciseID<10 {\input{build/Q3}}  \fi% 3
\ifnum\theexerciseID<10 {\input{build/Q4}}  \fi% 4
\ifnum\theexerciseID<10 {\input{build/Q5}}  \fi% 5
\ifnum\theexerciseID<10 {\input{build/Q6}}  \fi% 6
\ifnum\theexerciseID<10 {\input{build/Q7}}  \fi% 7
\ifnum\theexerciseID<10 {\input{build/Q8}}  \fi% 8
\ifnum\theexerciseID<10 {\input{build/Q9}}  \fi% 9
\ifnum\theexerciseID<10 {\input{build/Q10}} \fi%10
\ifnum\theexerciseID<10 {\input{build/Q11}} \fi%11
\ifnum\theexerciseID<10 {\input{build/Q12}} \fi%12
\ifnum\theexerciseID<10 {\input{build/Q13}} \fi%13
\ifnum\theexerciseID<10 {\input{build/Q14}} \fi%14
\ifnum\theexerciseID<10 {\input{build/Q15}} \fi%15
\ifnum\theexerciseID<10 {\input{build/Q16}} \fi%16
\ifnum\theexerciseID<10 {\input{build/Q17}} \fi%17
\ifnum\theexerciseID<10 {\input{build/Q18}} \fi%18
\ifnum\theexerciseID<10 {\input{build/Q19}} \fi%19
\ifnum\theexerciseID<10 {\input{build/Q20}} \fi%20


\end{document}
